\subsection{Robustness Checks}

To make sure that the reported results from the model and the following interpretation of the mechanisms driving segregation and inequalities are robust, I varied several of the arbitrary choices I had to make when implementing the model in reduced simulation experiments. Table \ref{tab:parameters} gives an overview of all parameters and which of them have been varied in which robustness check. I also implemented a version in NetLogo to make sure there are no major programming errors. The full analyses of all robustness checks can be found in the online supplement. All in all, the presented results are robust to arbitrary parameter choices and choice of programming language. However, some of the robustness checks reveal further insights about the model that I briefly want to highlight here.

In my motivation and review of existing theories, I made clear that I believe a model of housing markets should also be able to explain gentrification after a demand shock. While the main analysis suggests a dynamic process that mimics many features of gentrification, I wanted to explicitly test how my model reacts to a demand shock after converging to equilibrium. To do so, I ran the most likely parameter combination until equilibrium and simulated an increase/decrease in income inequality and/or an increase/decrease in population (for details, see online supplement robustness check C). Specific neighborhoods do not ascend or descend in this simulation experiment; instead, the shocks are absorbed equally throughout the city. I believe the reason for the lack of gentrification is that the model is in perfect equilibrium. Real cities have much lower levels of segregation than observed in my simulations, which means there are no mixed neighborhoods at risk of gentrification. The model abstracts from market frictions and path dependencies, which create pockets with more or less potential for redevelopment after such a demand shock. This is certainly a limitation of this model.

I also varied the inequality in income (robustness check E) between a Gini index of 0.26 and 0.42, in addition to the original level of 0.32. This analysis reveals two inconsistencies between my model and prior empirical research. First, different levels of income inequality do not lead to different levels of income segregation. Because every household will move to achieve even a marginal improvement in residential satisfaction, cities end up almost perfectly segregated no matter what. \citet{yavasDissectingIncomeSegregation2019} already made the same observation for utility-maximization models and therefore proposed a bounded rationality model that can explain the link between inequality and segregation. In a bounded rationality framework, making a segregating move becomes more likely as the absolute differences between neighborhoods increase. Second, in contrast to the findings of \citet{reardonIncomeInequalityIncome2011}, among others, income segregation by income quantile does not follow a U-shape in my data. While the lower quantiles of the income distribution tend to be very segregated, the richest households are not the most segregated from the rest of the income distribution but actually the least segregated. This could again be a consequence of the behavioral rule, as \citet{yavasDissectingIncomeSegregation2019} reproduces this result, or it could be due to the preference function: I assume that every household has the same preferences, whereas rich households might have a stronger desire to separate themselves from others. Together, these shortcomings indicate that future models should pay closer attention to their assumptions regarding human decision-making in residential contexts.
